\subsection{RQ4: What does not pay off}
\label{sec:rq4}
The previous results describe the regime where planning wins. Equally important for practitioners is where it does not, because a planner that is enabled by default must lose little there. We group the evidence by mechanism.

\textbf{Parallel workers for small tasks.} Cone reduction leaves little for workers to accelerate. On the held-out 12-qubit, depth-four QML probe (24 cases with fixed execution graphs and 1,152 validated collections), four workers reduced the summed collection time by 12\%, but per case they were 1.52 to 1.64 times slower than serial execution and won only 6 to 8 of the 24 cases. After charging one pool start-up and shutdown, the whole probe took 0.998 times the serial time, and two workers cost 1.32 times. Workers paid off only on wide independent components, with 3.18 to 3.28 times warm speedup for eight 16-qubit components on four workers (Appendix~\ref{app:parallel}). Planning and parallelism compete for the same slack: once the planner has shrunk the tasks, there is little left to distribute, which is also why we do not consider distributed execution part of the contribution.

\textbf{Exact cutting at small widths.} Joint cut selection does reduce unique fragment experiments, from 148 to 67 on a QAOA path, and repeated collections repay its search after 18 to 21 calls (Appendix~\ref{app:earlier}). None of these comparisons beats dense execution while the dense state fits, though: ten dense collections of a six-qubit path took 1.06\,ms against 389\,ms for the joint cut plan. With $10^k$ exact branch combinations for $k$ cuts, cutting is a fallback for cones that exceed the payload limit by a few qubits rather than an accelerator.

\textbf{ZX rewrites.} In a 12-qubit pilot with 64 Pauli terms, verified ZX rewriting spent 199\,ms of search to produce the same 11 groups as no rewriting, and in the cutting fixtures both policies selected original-circuit variants. We have not isolated a positive contribution of rewriting to plan cost.

\textbf{Quantum hardware.} Plans are ordinary circuits and can be submitted to hardware, but the trade-off changes there, because a QPU returns every same-basis output of a circuit from one set of shots. On IQM Garnet, with 20 same-state $Z$ outputs of a 20-qubit hardware-efficient circuit and 100 shots per circuit, the full circuit used one circuit and 0.70\,s of station execution time, while grouped cones needed 10 circuits and 1.06\,s, and separate cones 20 circuits and 1.40\,s. Cones can only help when few outputs are requested or when they need different measurement bases. We tested that case with $X$ and $Z$ on two wires of the same circuit, two seeds, 300 shots per output in every plan and a randomized plan order. A fair accuracy comparison needs every plan to read its outputs on the same physical qubits. Garnet's coupling graph has no 20-qubit Hamiltonian path, but it has a 19-qubit path that starts at the two output qubits, so we mapped the logical wires of every plan onto it. Cone circuits then need no routing, and only the full circuit routes its last wire, far from the outputs. With this placement, grouped cones ran in 0.56 times the full circuit's execution window at equal shots on both seeds, and in 0.69 to 0.85 times the window of separate cones with half their shots, using 14 two-qubit gates per circuit against 88. No accuracy difference between plans was resolved: the mean absolute error against exact values was 0.054 to 0.062, and every paired difference had a bootstrap interval spanning zero. An earlier run that moved the outputs into place with terminal SWAPs had found grouped cones less accurate than separate ones, and the repeat traces that deficit to their six extra SWAPs. Earlier two-output pilots with automatic and identity layouts pointed the same way but confounded placement. Two seeds on one device with sequential jobs say nothing about queues or scheduling, so the hardware evidence remains exploratory. IonQ's noisy simulator shows the same pattern in processing time (17\,s for full 20-qubit jobs against 1\,s for width-six cone jobs), which is not a hardware measurement.

\begin{table}[!t]
\caption{Trained Classifiers on One CPU Thread: Median Time over Direct Qulacs. Iris Rows Time Complete 30-Step Training Runs, Replayed Rows Time Inference Batches of 64 Examples Including Setup. F and G Denote the Full and Grouped Plans Chosen by \texttt{auto}.}
\label{tab:trained}
\centering
\setlength{\tabcolsep}{3pt}
\footnotesize
\resizebox{\columnwidth}{!}{\begin{tabular}{lrlrr}
\toprule
Models & $N$ & Auto choice & Auto/direct & Grouped/direct\\
\midrule
Iris, brickwork, 4 qubits & 12 & F, G & 1.77 & 1.64\\
Iris, brickwork, 8 qubits & 12 & F & 1.34 & 1.42\\
Iris, ring, 4 qubits & 12 & F & 1.81 & 2.19\\
Iris, ring, 8 qubits & 12 & F & 1.37 & 3.54\\
\midrule
MNIST1D, 8 qubits & 5 & F & 1.23 & 5.59\\
EEG, 8 qubits & 5 & F & 1.27 & 5.66\\
PneumoniaMNIST, 8 qubits & 5 & F & 1.24 & 5.49\\
\bottomrule
\end{tabular}
}
\end{table}

\textbf{Small trained classifiers.} Trained models at the widths where we could train them do not benefit (Table~\ref{tab:trained}). For each of the 48 Iris classifiers, direct Qulacs was faster than every planned method, by 1.3\% in the closest case, and \texttt{auto} cost 1.34 to 1.81 times as much over a complete training run. Each step changes the parameters and therefore builds a new plan, and the circuits are too small for cones to repay it. The brickwork ansatz keeps cones local, four of eight qubits at depth two, yet grouped cones still cost 1.42 to 1.64 times direct execution. The 15 replayed classifiers saturate their cones with four re-uploads of a CX ring. \texttt{auto} kept the dense plan for all of them, within 2\% of the forced dense plan, and still cost 1.23 to 1.27 times direct Qulacs, because \textsc{LazyQC} lowers native circuits at every collection. Grouped cones, which ignore the saturation, cost 5.5 to 5.7 times. Replaying these models also checks the circuit translation against a different simulator and code base: all 3,300 saved test predictions were reproduced exactly. The planning cache cannot help such inference either, since new data change the bindings that its keys include.

\textbf{Approximation without structural change.} The opt-in sensitivity mechanism of Appendix~\ref{app:sensitivity} pays off only when an admissible substitution changes the structure of the plan. On constructed 16-qubit circuits with weak bridges it halves the maximum width and gives 2.89 to 3.57 times speedup over eight fully charged collections, while a connected control and all eight-qubit fixtures regress once profiling is charged.
